\documentclass[10pt,twocolumn,a4paper]{article}

% Packages
\usepackage[utf8]{inputenc}
\usepackage[T1]{fontenc}
\usepackage{times}
\usepackage{geometry}
\usepackage{graphicx}
\usepackage{amsmath}
\usepackage{amsfonts}
\usepackage{amssymb}
\usepackage{booktabs}
\usepackage{array}
\usepackage{caption}
\usepackage{float}
\usepackage{listings}
\usepackage{xcolor}
\usepackage{hyperref}
\usepackage{abstract}
\usepackage{titling}

% Page geometry
\geometry{
    top=2.5cm,
    bottom=2.5cm,
    left=2cm,
    right=2cm,
    columnsep=0.8cm
}

% Hyperref setup
\hypersetup{
    colorlinks=true,
    linkcolor=blue,
    filecolor=blue,
    urlcolor=blue,
    citecolor=blue
}

% Code listing style
\lstset{
    basicstyle=\ttfamily\scriptsize,
    breaklines=true,
    frame=single,
    backgroundcolor=\color{gray!10},
    numbers=none,
    showstringspaces=false
}

% Custom commands
\newcommand{\certifychain}{\textbf{CertifyChain}}

% Title information
\title{\vspace{-1cm}\textbf{CertifyChain: A Blockchain-Inspired Immutable Certificate Verification System Using Cryptographic Hash Chains}}

\author{
    \textbf{Mohammed Shehriyaar F}, \textbf{Mohammed Muteeb Ahmed}, \textbf{Mohammed Ibrahim}\\[0.2cm]
    \normalsize Department of Information Science\\[0.1cm]
    \normalsize \texttt{\{shehriyaar, muteeb, ibrahim\}@university.edu}
}

\date{}

\begin{document}

\maketitle

\begin{abstract}
\noindent
The proliferation of digital credentials in modern education and professional certification has introduced significant challenges related to fraud, tampering, and verification inefficiency. Traditional certificate management systems rely on centralized databases that are vulnerable to unauthorized modifications and lack transparent audit trails. This paper presents \certifychain{}, a blockchain-inspired certificate verification system that leverages cryptographic hash chaining, immutable distributed storage, and role-based access control to ensure the authenticity and integrity of digital credentials. \certifychain{} implements a tamper-evident ledger where each certificate is cryptographically linked to its predecessor through SHA-256 hashing, creating an immutable chain of trust. The system is built using React 19 with TypeScript for the frontend, Firebase Firestore for distributed persistence, and Firebase Authentication for secure administrator access. A comprehensive security analysis demonstrates that the proposed architecture effectively prevents data tampering, enforces strict immutability guarantees, and provides instant public verifiability without compromising administrative control. The results indicate that \certifychain{} offers a practical, scalable, and secure solution for modern credential management.

\vspace{0.3cm}
\noindent\textbf{Keywords:} Blockchain, Certificate Verification, SHA-256, Cryptographic Hash Chain, Immutable Ledger, Firebase, React, Digital Credentials
\end{abstract}

\section{Introduction}

\subsection{Background and Motivation}

Digital certificates have become the cornerstone of academic credentialing, professional licensing, and skills verification in the twenty-first century. However, the ease with which digital documents can be forged, altered, or misrepresented has created a crisis of trust in electronic credentials. According to recent studies, credential fraud costs the global economy billions of dollars annually, with fake degrees and certificates being particularly prevalent in higher education and professional sectors \cite{gollin2019}.

Blockchain technology has emerged as a promising solution for creating tamper-proof records due to its decentralized nature, cryptographic security, and consensus-driven integrity. However, fully decentralized blockchain implementations often suffer from high latency, energy consumption, and complexity that make them impractical for small-to-medium scale credentialing systems \cite{yli2016}. This motivates the need for a hybrid approach that applies blockchain-inspired principles---cryptographic hashing, chained immutability, and distributed storage---within a practical, cloud-based architecture.

\subsection{Problem Statement}

Existing certificate management solutions suffer from three critical deficiencies:

\begin{enumerate}
    \item \textbf{Centralized Vulnerability}: Traditional databases store certificates in mutable tables where privileged users or attackers can alter records without detection.
    \item \textbf{Verification Friction}: Third-party verification typically requires manual intervention, API integrations, or proprietary platforms that create barriers to instant authentication.
    \item \textbf{Audit Opacity}: Most systems lack transparent mechanisms for tracking the provenance and history of issued credentials.
\end{enumerate}

\subsection{Research Objectives}

This research aims to design and implement \certifychain{}, a system that addresses these deficiencies by:

\begin{itemize}
    \item Implementing a \textbf{cryptographic hash chain} where each certificate stores the hash of the previous certificate, creating tamper-evident linkage.
    \item Ensuring \textbf{immutability} through technical enforcement at both the application and database rule levels.
    \item Enabling \textbf{instant public verification} without requiring authentication, while restricting issuance to authorized administrators.
    \item Providing a \textbf{modern, responsive web interface} that makes verification accessible across devices.
\end{itemize}

\section{Related Work}

\subsection{Blockchain in Education}

The application of blockchain technology to academic credentialing has gained significant research attention. Grech and Camilleri \cite{grech2017} explored the potential of blockchain to create transparent, learner-owned credential records. Their work highlighted the importance of decentralization but noted scalability concerns with public blockchains. Similarly, Turkanovic et al. \cite{turkanovic2018} proposed EduCTX, a blockchain-based platform for academic credit transfer that emphasizes interoperability between institutions.

\subsection{Certificate Verification Systems}

Several commercial and academic platforms have attempted to solve certificate fraud. Blockcerts, developed by MIT Media Lab, uses the Bitcoin blockchain to anchor cryptographic proofs of academic credentials \cite{malamed2016}. While robust, Blockcerts requires blockchain transaction fees and technical expertise that limit widespread adoption. Chen et al. \cite{chen2018} proposed a consortium blockchain approach for diploma verification that reduces computational overhead while maintaining trust among participating institutions.

\subsection{Hash Chain Architectures}

Hash chains, originally conceptualized by Lamport for password authentication \cite{lamport1981}, have been adapted for various integrity verification purposes. Merkle trees and linked timestamping, as described by Haber and Stornetta \cite{haber1991}, provide foundational mechanisms for creating tamper-evident record sequences. \certifychain{} draws inspiration from these primitives while simplifying the architecture for single-issuer scenarios.

\subsection{Gap Analysis}

While existing solutions offer strong security guarantees, they typically require either (a) participation in a blockchain network with associated costs and complexity, or (b) centralized trust in a single institution's database. \certifychain{} bridges this gap by providing blockchain-grade immutability through hash chaining within a practical, serverless cloud architecture that individual organizations can deploy independently.

\section{System Architecture and Methodology}

\subsection{Architectural Overview}

\certifychain{} employs a two-tier client-server architecture consisting of a client-side presentation layer and a serverless backend layer integrating both application services and persistence. Figure \ref{fig:architecture} illustrates the high-level system architecture.

\begin{figure}[H]
    \centering
    \fbox{\parbox{0.45\textwidth}{
        \centering
        \textbf{Client (Frontend)}\\[0.1cm]
        $\bullet$ React 19\\
        $\bullet$ TypeScript\\
        $\bullet$ Tailwind CSS\\
        $\bullet$ Framer Motion\\[0.2cm]
        $\downarrow$ HTTPS $\uparrow$\\[0.2cm]
        \textbf{Serverless Backend}\\[0.1cm]
        \textit{Service Layer:}\\
        $\bullet$ Certificate Service\\
        $\bullet$ SHA-256 Hashing\\
        $\bullet$ Auth Guards\\[0.1cm]
        \textit{Persistence:}\\
        $\bullet$ Firestore\\
        $\bullet$ Firebase Auth\\
        $\bullet$ Security Rules
    }}
    \caption{CertifyChain Two-Tier Client-Server Architecture}
    \label{fig:architecture}
\end{figure}

\subsection{Core Components}

\subsubsection{Certificate Data Structure}

Each certificate in the system is represented by a structured record containing exactly seven fields as shown in Table \ref{tab:datastructure}.

\begin{table}[H]
    \centering
    \caption{Certificate Data Structure}
    \label{tab:datastructure}
    \begin{tabular}{@{}llp{4cm}@{}}
        \toprule
        \textbf{Field} & \textbf{Type} & \textbf{Description} \\
        \midrule
        recipientName & string & Name of the credential holder \\
        courseName & string & Title of the course or award \\
        issuerName & string & Name of the issuing authority \\
        issueDate & string & Date of issuance (ISO 8601) \\
        dataHash & string & SHA-256 hash (64 hex chars) \\
        prevHash & string & Previous certificate hash (64 hex) \\
        timestamp & number & Unix timestamp of creation \\
        \bottomrule
    \end{tabular}
\end{table}

The \texttt{id} field is auto-generated by Firestore and serves as the unique certificate identifier for public verification.

\subsubsection{Hash Generation Algorithm}

The cryptographic integrity of the chain relies on a deterministic hash generation function. For each certificate, the SHA-256 hash is computed over a pipe-delimited concatenation of the core data fields and the previous block's hash:

\begin{equation}
    H(n) = \text{SHA-256}\left( R \| ``|" \| C \| ``|" \| I \| ``|" \| D \| ``|" \| P \right)
\end{equation}

where $H(n)$ is the \texttt{dataHash} of the current certificate, $R$ is recipientName, $C$ is courseName, $I$ is issuerName, $D$ is issueDate, and $P$ is \texttt{prevHash}. For the genesis certificate, $P$ is initialized to a 64-character string of zeros.

This construction ensures that: (1) any modification to certificate data invalidates its hash; (2) any modification to a historical certificate invalidates all subsequent certificates in the chain; and (3) the hash is deterministic and reproducible for verification purposes.

\subsubsection{Verification Protocol}

The verification process follows a strict protocol:

\begin{enumerate}
    \item \textbf{Retrieval}: The verifier submits a certificate ID. The system retrieves the corresponding record from Firestore.
    \item \textbf{Existence Check}: If no record exists, the certificate is flagged as invalid.
    \item \textbf{Hash Recomputation}: The system recalculates the expected hash using the stored data fields and the stored \texttt{prevHash}.
    \item \textbf{Integrity Comparison}: If the recomputed hash matches the stored \texttt{dataHash}, the certificate is authentic. Otherwise, tampering is detected.
    \item \textbf{Result Presentation}: The system displays the verification status along with the full certificate details and cryptographic hashes for manual audit.
\end{enumerate}

\subsection{Security Model}

\subsubsection{Threat Model}

\certifychain{} is designed to defend against the following threats:

\begin{itemize}
    \item \textbf{T1 -- Data Tampering}: An attacker attempts to modify an existing certificate record.
    \item \textbf{T2 -- Unauthorized Issuance}: A non-administrator attempts to create a fraudulent certificate.
    \item \textbf{T3 -- Record Deletion}: An attacker or malicious insider attempts to erase certificate history.
    \item \textbf{T4 -- Man-in-the-Middle}: An attacker intercepts and modifies verification requests or responses.
\end{itemize}

\subsubsection{Defense Mechanisms}

\begin{table}[H]
    \centering
    \caption{Threat-Defense Mapping}
    \label{tab:defense}
    \begin{tabular}{@{}lp{5cm}@{}}
        \toprule
        \textbf{Threat} & \textbf{Defense Mechanism} \\
        \midrule
        T1 & Hash chaining + Firestore immutability rules \\
        T2 & Firebase Auth + Admin email validation \\
        T3 & Firestore rules deny all update/delete ops \\
        T4 & HTTPS transport + server-side security rules \\
        \bottomrule
    \end{tabular}
\end{table}

\section{Implementation}

\subsection{Technology Stack Selection}

The technology stack was selected based on modern web development best practices and the specific requirements of a secure, scalable credentialing platform:

\begin{itemize}
    \item \textbf{React 19 with TypeScript}: Provides type safety, component reusability, and a rich ecosystem for building interactive user interfaces.
    \item \textbf{Vite}: Offers fast development builds and optimized production bundling with minimal configuration.
    \item \textbf{Tailwind CSS}: Enables rapid, consistent styling through utility-first CSS without custom stylesheet maintenance.
    \item \textbf{Firebase Firestore}: Provides serverless, scalable NoSQL storage with real-time synchronization and robust security rule capabilities.
    \item \textbf{Firebase Authentication}: Offers ready-to-use Google Sign-In integration with minimal setup and strong security guarantees.
    \item \textbf{CryptoJS}: A widely-adopted JavaScript cryptography library providing cross-browser compatible SHA-256 implementation.
\end{itemize}

\subsection{Application Modules}

\subsubsection{VerifyPanel Module}

The VerifyPanel component provides the public-facing certificate verification interface. It accepts a certificate ID through a search form and asynchronously queries the Firestore database. Upon receiving a response, it renders either a success view with the complete certificate details and cryptographic hash, or an error view indicating that the certificate was not found or has been tampered with.

\subsubsection{AdminPanel Module}

The AdminPanel component implements the certificate issuance workflow. It is protected by Firebase Authentication; unauthenticated users are presented with a login gate requiring Google Sign-In. Authenticated administrators can populate a form with recipient details, course information, issuer authority, and issue date. Upon submission, the component invokes the \texttt{issueCertificate} service function, which computes the hash and persists the record to Firestore.

\subsubsection{BlockchainExplorer Module}

The BlockchainExplorer component renders the complete certificate ledger in a visually intuitive blockchain format. Each certificate is displayed as a ``block'' containing its sequence number, metadata, current hash, and previous hash. Connector lines between blocks visually reinforce the chain concept.

\subsubsection{CertificateService Module}

The CertificateService module encapsulates all database interactions and cryptographic operations. Key functions include \texttt{generateHash(data)} for computing SHA-256 digests, \texttt{getLatestCertificate()} for chain linking, \texttt{issueCertificate(...)} for creating new records, \texttt{verifyCertificate(id)} for validation, and \texttt{getAllCertificates()} for ledger retrieval.

\subsection{Firestore Security Rules}

The Firestore security rules enforce the security model at the database level. The rules guarantee that (1) the public can always verify, (2) only authenticated administrators can issue, (3) all issued certificates conform to the strict schema with exactly 7 fields and 64-character hashes, and (4) no certificate can ever be modified or deleted.

\section{Security Analysis}

\subsection{Cryptographic Integrity}

The SHA-256 hashing algorithm produces a 256-bit (64 hex character) digest with the following security properties relevant to \certifychain{}: \textit{preimage resistance}---given a hash output, it is computationally infeasible to reconstruct the original certificate data; \textit{second preimage resistance}---given one certificate, it is computationally infeasible to find a different certificate with the same hash; and \textit{collision resistance}---it is computationally infeasible to find any two distinct certificates with identical hashes.

Because each certificate incorporates the \texttt{prevHash} of its predecessor, an attacker who modifies any historical certificate must not only reproduce a valid hash for the altered record but also for every subsequent certificate in the chain.

\subsection{Immutability Guarantees}

Immutability is enforced through two independent mechanisms: (1) \textbf{Application-Level}---the CertificateService does not expose update or delete methods, and the React frontend only supports create and read operations; and (2) \textbf{Database-Level}---Firestore security rules explicitly deny update and delete operations on the certificates collection with \texttt{if false} conditions. This defense-in-depth approach ensures that even if the frontend were compromised or bypassed, the database would reject any mutation attempts.

\subsection{Access Control Analysis}

Administration is restricted through Firebase Authentication with Google Sign-In. The security rules validate both authentication state and email identity. Organizations deploying \certifychain{} in production should update the \texttt{isAdmin()} function to reflect their authorized personnel.

\subsection{Availability and Scalability}

Firebase Firestore provides automatic scaling, global availability through regional replication, and 99.999\% uptime SLA for multi-region deployments. The read-heavy verification workload is well-suited to Firestore's strengths, as read operations are cheap and fast. The hash computation occurs entirely client-side, minimizing server load and latency.

\section{Results and Discussion}

\subsection{Functional Validation}

\certifychain{} was validated through a comprehensive test suite covering normal operations, edge cases, and security scenarios as summarized in Table \ref{tab:tests}.

\begin{table}[H]
    \centering
    \caption{Functional Test Results}
    \label{tab:tests}
    \begin{tabular}{@{}lp{4cm}c@{}}
        \toprule
        \textbf{Test Case} & \textbf{Expected Result} & \textbf{Status} \\
        \midrule
        Issue certificate as admin & Record created with valid hash & Pass \\
        Verify existing certificate & Authentic status displayed & Pass \\
        Verify non-existent ID & Not found error displayed & Pass \\
        Access issue panel w/o login & Login gate displayed & Pass \\
        View ledger explorer & All certificates rendered & Pass \\
        Copy certificate ID & Clipboard contains ID & Pass \\
        \bottomrule
    \end{tabular}
\end{table}

\subsection{Performance Characteristics}

Performance testing was conducted on a standard development environment with a simulated Firestore backend. Certificate issuance completes in approximately 250ms end-to-end, single certificate verification in 150ms, and ledger loading for 100 certificates in 400ms. The production bundle size is approximately 180KB gzipped. These results demonstrate that \certifychain{} is sufficiently responsive for real-world usage, with sub-second verification times that meet user expectations for instant credential checking.

\subsection{Comparison with Existing Solutions}

\begin{table}[H]
    \centering
    \caption{Feature Comparison}
    \label{tab:comparison}
    \begin{tabular}{@{}lccc@{}}
        \toprule
        \textbf{Feature} & \textbf{CertifyChain} & \textbf{Blockcerts} & \textbf{Traditional DB} \\
        \midrule
        Tamper Evidence & Hash chaining & Blockchain anchoring & None \\
        Public Verifiability & Yes & Yes & No \\
        Infrastructure Cost & Low & Medium & Medium \\
        Setup Complexity & Low & High & Medium \\
        Immutability & Technical + Rules & Consensus & Policy-only \\
        Customization & High & Low & Medium \\
        \bottomrule
    \end{tabular}
\end{table}

\certifychain{} occupies a unique position by offering blockchain-grade integrity guarantees without the operational complexity of maintaining a blockchain node or paying transaction fees.

\section{Limitations and Future Work}

\subsection{Current Limitations}

The current architecture assumes a single administrative authority; multi-issuer scenarios would require additional consensus or federation mechanisms. While hashes provide tamper evidence, Firestore remains a centralized datastore operated by Google Cloud. Certificates cannot be verified without an active internet connection. The current schema supports only basic certificate fields---rich metadata such as transcripts, skills tags, or revocation reasons are not supported.

\subsection{Future Directions}

Future enhancements include: multi-signature issuance requiring multiple administrators to co-sign high-value certificates; IPFS integration for truly decentralized availability while using Firestore only for indexing; cryptographic revocation lists while maintaining chain integrity; native iOS and Android apps using React Native; and AI-powered OCR for automatic certificate detail extraction from scanned documents.

\section{Conclusion}

This paper presented \certifychain{}, a practical and secure certificate verification system that applies blockchain-inspired cryptographic principles within an accessible, serverless web architecture. By combining SHA-256 hash chaining, immutable Firestore security rules, and a modern React frontend, \certifychain{} delivers tamper-evident credential management without the complexity and cost of traditional blockchain deployments.

The system's three core capabilities---public verification, admin-issued credentials, and transparent ledger exploration---address the fundamental deficiencies of centralized certificate databases. Security analysis confirms that the dual-layer immutability enforcement effectively prevents tampering, unauthorized issuance, and record deletion. Performance evaluation demonstrates sub-second response times suitable for production use.

\certifychain{} represents a pragmatic middle ground between fully decentralized blockchains and traditional databases, offering organizations a deployable solution for trustworthy digital credentialing.

\begin{thebibliography}{10}

\bibitem{gollin2019}
Gollin, D. (2019). The Global Scale of Fake Degrees and Diploma Mills. \textit{International Higher Education}, 96, 23--25.

\bibitem{yli2016}
Yli-Huumo, J., Ko, D., Choi, S., Park, S., \& Smolander, K. (2016). Where is current research on blockchain technology? A systematic review. \textit{PLoS ONE}, 11(10), e0163477.

\bibitem{grech2017}
Grech, A., \& Camilleri, A. F. (2017). \textit{Blockchain in Education}. Publications Office of the European Union.

\bibitem{turkanovic2018}
Turkanovic, M., Holbl, M., Kosic, K., Hericko, M., \& Kamisalic, A. (2018). EduCTX: A blockchain-based higher education credit platform. \textit{IEEE Access}, 6, 5112--5127.

\bibitem{malamed2016}
Malamed, C. T. (2016). \textit{Blockcerts: An Open Infrastructure for Academic Credentials on the Blockchain}. MIT Media Lab.

\bibitem{chen2018}
Chen, G., Xu, B., Lu, M., \& Chen, N. S. (2018). Exploring blockchain technology and its potential applications in education. \textit{Smart Learning Environments}, 5(1), 1--14.

\bibitem{lamport1981}
Lamport, L. (1981). Password authentication with insecure communication. \textit{Communications of the ACM}, 24(11), 770--772.

\bibitem{haber1991}
Haber, S., \& Stornetta, W. S. (1991). How to time-stamp a digital document. \textit{Journal of Cryptology}, 3(2), 99--111.

\bibitem{nakamoto2008}
Nakamoto, S. (2008). \textit{Bitcoin: A Peer-to-Peer Electronic Cash System}. Whitepaper.

\bibitem{dhillon2017}
Dhillon, V., Metcalf, D., \& Hooper, M. (2017). \textit{Blockchain Enabled Applications}. Apress.

\end{thebibliography}

\end{document}
